% Pharo poster -- A1 portrait.  Compile TWICE with pdflatex, lualatex or xelatex (TikZ overlay needs 2 runs)
% Figures expected in figures/: pharo-logo.pdf, pharo-beacon.pdf
% Also uses figures/inria-transparent.pdf: newINRIALogo_En.pdf without its white
% background and with the grey "informatics mathematics" text turned white.
\documentclass{article}
\usepackage[paperwidth=594mm,paperheight=841mm,margin=0mm]{geometry}
\usepackage{iftex}
\ifPDFTeX            % pdflatex: Type1 Fira Sans
  \usepackage[T1]{fontenc}
  \usepackage{lmodern}
  \usepackage[sfdefault,light,semibold]{FiraSans}
  \newcommand{\heavy}{\fontseries{b}\selectfont}
\else                % lualatex / xelatex: OpenType Fira Sans
  \usepackage{fontspec}
  \setmainfont{Fira Sans}[
    UprightFont=*-Light, BoldFont=*-SemiBold, ItalicFont=*-LightItalic, BoldItalicFont=*-SemiBoldItalic]
  \newfontfamily\heavy{Fira Sans}[UprightFont=*-Bold, ItalicFont=*-BoldItalic]
\fi
\usepackage{tikz}
\usetikzlibrary{calc}
\usepackage{fontawesome5}
\pagestyle{empty}
\hyphenpenalty=10000\exhyphenpenalty=10000
\setlength{\parindent}{0pt}

% ---------- colours ----------
\definecolor{cream}{HTML}{F7F4EE}
\definecolor{greybg}{HTML}{DCDEE6}
\definecolor{blueA}{HTML}{4A78B5}
\definecolor{blueB}{HTML}{2E5089}
\definecolor{stripe}{HTML}{3E68A6}
\definecolor{ink}{HTML}{1F2A3A}
\definecolor{osigreen}{HTML}{3DA639}

% ---------- helpers ----------
\newcommand{\fs}[1]{\fontsize{#1}{1.2\dimexpr#1pt\relax}\selectfont}
\newcommand{\num}[1]{{\heavy\fs{52}#1}}                 % big number
\newcommand{\stat}[2]{\num{#1}\,\,#2\par\vspace{4mm}}   % number + label
\newcommand{\mstat}[2]{{\heavy\fs{42}#1}\,\,{\fs{28}#2}\par\vspace{2mm}} % compact stat for the MOOC columns
\newcommand{\head}[1]{{\heavy\fs{62}#1}}                 % section heading
\newcommand{\web}{\faGlobe\,\,}                          % marks a web address
\newcommand{\url}[1]{{\fs{32}\web#1}}
\newcommand{\headurl}[2]{\head{#1}\\[2mm]\url{#2}}   % heading with its web address right below, left-aligned

\begin{document}
\fs{31}%
\begin{tikzpicture}[remember picture, overlay,
    x=\paperwidth, y=-\paperheight, shift={(current page.north west)}]

  % ----- key points of the geometry (x,y as fractions; y goes down) -----
  \coordinate (TL)  at (0,0);     \coordinate (TR) at (1,0);
  \coordinate (BL)  at (0,1);     \coordinate (BR) at (1,1);
  \coordinate (S1)  at (0.63,0);  \coordinate (S2) at (0.69,0);     % stripe top
  \coordinate (C1)  at (0.43,0.25); \coordinate (C2) at (0.50,0.25); % stripe foot
  \coordinate (M1)  at (0.69,0.635); \coordinate (M2) at (0.72,0.62); % middle joint
  \coordinate (G2)  at (1,0.60);                                     % grey bottom right
  \coordinate (B1)  at (0.63,1);                                     % blue bottom right

  % ----- background panels -----
  \fill[cream] (TL) rectangle (BR);
  % grey lighthouse panel
  \fill[greybg] (S2) -- (TR) -- (G2) -- (M2) -- (C2) -- cycle;
  % blue diagonal stripe carrying pharo.org
  \fill[stripe] (S1) -- (S2) -- (C2) -- (C1) -- cycle;
  \node[rotate=61, text=white, font=\fs{48}, anchor=center]
       at (0.607,0.075) {pharo.org};
  % big blue panel
  \shade[top color=blueA, bottom color=blueB]
       (0,0.25) -- (C1) -- (M1) -- (B1) -- (BL) -- cycle;
  % darker diagonal shading inside the blue panel
  \fill[black, opacity=0.13]
       (0.18,0.62) -- (0.55,0.45) -- (M1) -- (B1) -- (0.12,1) -- cycle;

  % ----- large lighthouse, clipped to the grey panel -----
  \begin{scope}
    \clip (S2) -- (TR) -- (G2) -- (M2) -- (C2) -- cycle;
    \node at (0.825,0.325) {\includegraphics[height=0.64\paperheight]{figures/pharo-beacon.pdf}};
  \end{scope}

  % ----- header: logo + tagline + intro -----
  \node[anchor=north west, inner sep=0] (logo) at (0.07,0.015)
       {\includegraphics[width=0.42\paperwidth]{figures/pharo-logo.pdf}};
  \node[anchor=north west, text=blueB, font=\fs{36}] at (0.13,0.118)
       {The live programming experience};

  \node[anchor=north west, text width=0.40\paperwidth, text=ink, align=left, font=\fs{27}]
       at (0.08,0.145) {%
    Pharo is a pure open-source reflective object-oriented programming
    language and a powerful environment, focused on simplicity and
    immediate feedback. Fully written in itself.\par\vspace{6mm}
    Pharo offers a unique immersive programming experience.
    Amazing debugging experience.};

  % ----- blue panel content -----
  \begin{scope}[every node/.style={text=white, anchor=north west, align=left}]

  % key figures
  \node[text width=0.48\paperwidth] at (0.11,0.27) {%
    Full syntax on half of a postcard\par\vspace{4mm}
    \stat{740}{internal packages}
    More than \num{3000}\,\,external projects\par\vspace{4mm}
    \stat{30}{Universities use Pharo to teach OOP}
    \stat{15}{international research groups use Pharo}
    \stat{98}{committers}};

  % industrial consortium
  \node at (0.075,0.417) {\fs{80}\faIndustry};
  \node at (0.165,0.410) {\headurl{Industrial Consortium}{consortium.pharo.org}};
  \node[text width=0.50\paperwidth] at (0.10,0.459) {%
    To structure the community\\
    To sustain the development of Pharo\par\vspace{1mm}
    \num{21}\,\,industrial members /\,\,\num{28}\,\,academic members};

  % association
  \node at (0.10,0.535) {\headurl{1 User Association}{association.pharo.org}};
  \node at (0.47,0.533) {\fs{100}\faUsers};

  % books
  \node at (0.10,0.603) {\headurl{29 open-source books}{books.pharo.org}};
  \node at (0.52,0.607) {\fs{100}\faBookOpen};

  % moocs: shared heading, then one column per mooc
  \node at (0.07,0.672) {\fs{110}\faGraduationCap};
  \node at (0.18,0.672) {\head{Two free MOOCs}};
  \node at (0.18,0.708) {\num{15 K}\,\,participants from 100+ countries};
  % column 1: Live Object Programming in Pharo
  \node[text width=0.25\paperwidth] at (0.075,0.752) {%
    {\heavy\fontsize{36}{39}\selectfont Live Object\\Programming in Pharo}\par\vspace{2mm}
    {\fs{24}\web mooc.pharo.org}\par\vspace{5mm}
    \mstat{93}{videos, 155 quiz}
    \mstat{6}{large exercises}
    \mstat{1}{full web application}
    \mstat{7}{challenges}
    {\fs{24}\itshape Dubbed EN \& FR\\Subtitles EN, FR, ES, JA}};
  % thin separator between the two columns
  \draw[white, opacity=0.5, line width=0.8mm] (0.345,0.755) -- (0.345,0.885);
  % column 2: Advanced Object-Oriented Design
  \node[text width=0.26\paperwidth] at (0.37,0.752) {%
    {\heavy\fontsize{36}{39}\selectfont Advanced Object-\\Oriented Design}\par\vspace{2mm}
    {\fs{24}\web advanced-design-mooc.pharo.org}\par\vspace{5mm}
    \mstat{10}{modules}
    \mstat{64}{lectures}
    \mstat{1}{exercise book}
    {\fs{28}Quizzes \& Open Badge}\par\vspace{3mm}
    {\fs{24}\itshape Concepts valid in any OO language}};
  \end{scope}

  % Inria logo, directly on the blue panel (transparent background)
  \node[anchor=center, inner sep=0] at (0.40,0.945)
    {\includegraphics[width=0.17\paperwidth]{figures/inria-transparent.pdf}};
  % QR code to pharo.org, left of the Inria logo (keeps its white quiet zone to stay scannable)
  \node[anchor=east, inner sep=0] at (0.285,0.945)
    {\includegraphics[width=0.08\paperwidth]{figures/qrcode-pharo.pdf}};

  % ----- technology stack (bottom right) -----
  \node[anchor=north east, align=right, text=ink] at (0.905,0.652)
       {\head{Technology}\\[-4mm]\head{Stack}};
  \node[anchor=north west, text=ink] at (0.913,0.660) {\fs{90}\faPuzzlePiece};

  \node[anchor=north west, text width=0.27\paperwidth, text=ink, align=left, font=\fs{25}]
       at (0.745,0.715) {%
    \textbf{AI/MCP}\\
    \textbf{Moose}\\
    \textbf{Bloc, Toplo, Roassal}\\
    \textbf{Metamodeling}\\
    \textbf{Web development:}\\
    \hspace*{6mm}- Seaside\\
    \hspace*{6mm}- Teapot, REST\\
    \textbf{Integration:} Jenkins, GitHub\\
    \textbf{Persistence:}\\
    \hspace*{6mm}- OO: GemStone\\
    \hspace*{6mm}- NoSQL: MongoDB, UnQLite,\\\hspace*{12mm}Tokyo, Redis\\
    \hspace*{6mm}- Relational: Postgres,\\\hspace*{12mm}MySQL, SQLite\\
    \textbf{Platforms:}\\
    \hspace*{6mm}- Linux, Windows, macOS\\
    \hspace*{6mm}- Mobile: iOS\\
    \hspace*{6mm}- Raspbian};

  % open source initiative badge
  \node[anchor=south, align=center, text=ink, font=\fs{20}] at (0.68,0.988) {%
    {\color{osigreen}\fs{95}\faOsi}\\[2mm]
    {\heavy open source}\\[-1mm] initiative\\[-1mm] MIT};
\end{tikzpicture}
\end{document}
